\documentclass[utf8, russian]{eskdtext}
\usepackage{paralist}
\usepackage{caption} %заголовки плавающих объектов
\captionsetup[table]{justification=centering} %установки для заголовков таблиц
% без определения что-то не работает
\newcommand{\No}{\textnumero}

% Fix for hyperref/eskdx conflict in modern TeXLive
\makeatletter
\@ifundefined{theHpoint}{}{\let\theHpoint\@undefined}
\@ifundefined{theHsubpoint}{}{\let\theHsubpoint\@undefined}
\@ifundefined{theHsubsubpoint}{}{\let\theHsubsubpoint\@undefined}
\makeatother

\usepackage{titlesec} 
\titleformat{\section}
{\normalsize\bfseries} 
{\thesection} {1em}{} 
\titleformat{\subsection}
{\normalsize\bfseries}
{\thesubsection} {1em}{} 
\titleformat{\subsubsection}
{\normalsize\bfseries}
{\thesubsubsection} {1em}{}




% для штампа
\ESKDtitle{{\small Название курсового проекта}}
\ESKDsignature{ТПЖА.090301.ХХХ ПЗ}
\ESKDgroup{\footnotesize Кафедра ЭВМ\\ИВТб-2301-06-01}
\ESKDauthor{Иванов}
\ESKDchecker{Петров}
%\ESKDnormContr{\resizebox{2.22cm}{\height}{Фамилия И.О.}}
%\ESKDapprovedBy{Фамилия И.О.}

% оступы от заголовков разделов и подразделов
%\ESKDsectSkip{section}{7mm}{7mm}
%\ESKDsectSkip{subsection}{5mm}{5mm}
%\ESKDsectSkip{subsubsection}{3mm}{3mm}

%%% Каждый раздел (\section) с новой страницы
\let\stdsection\section
\renewcommand \section{\newpage\stdsection}

\title{{\small Название курсового проекта}}

\usepackage{enumitem}
%\setlist[enumerate,itemize]{leftmargin=0pt,itemindent=2cm}
\renewcommand{\labelitemi}{\normalfont\bfseries{--}}
\setlist{nosep,
left=0pt,
labelindent=\parindent,leftmargin=*,
leftmargin=0pt,itemindent=\labelwidth
}
\setlist[enumerate,2]{
leftmargin=\parindent,itemindent=\labelwidth
}
\setenumerate{listparindent=\parindent}
\makeatletter
\AddEnumerateCounter{\asbuk}{\@asbuk}{м)}
\makeatother
\renewcommand{\labelenumi}{\asbuk{enumi})} 
\renewcommand{\labelenumii}{\arabic{enumii})}

%%% Выравнивание и переносы %%%

\sloppy                             % Избавляемся от переполнений

\clubpenalty=10000                  % Запрещаем разрыв страницы после первой строки абзаца

\widowpenalty=10000                 % Запрещаем разрыв страницы после последней строки абзаца
% подключение частей документа

%Для кода
\usepackage{listings}
\usepackage{verbatim}
%Для ссылок
\usepackage{hyperref}

\begin{document}
\maketitle
%%% Делаем содержание
\tableofcontents
%%% Введение пишется без цифры и добавляется в содержание
\section*{Введение}
\addcontentsline{toc}{section}{Введение}
Текст текст текст

\newpage

\section{Раздел 1}
Обзор предметной области и анализ аналогов
Разговорные и облачные роботы-компаньоны (например, AIBI)
Такие устройства в основном ориентированы на голосовое взаимодействие с пользователем. Они поддерживают диалог, выполняют голосовые команды и могут использовать внешние большие языковые модели (LLM) для обработки запросов.
Основным недостатком такого подхода является зависимость от внешних серверов и подключения к Интернету. Кроме того, подобные устройства обычно не ориентированы непосредственно на контроль концентрации пользователя во время работы.
Развлекательные и интерактивные роботы (например, Yaboo, Doll-E, Eilik)
Данная группа устройств ориентирована в первую очередь на эмоциональное и визуальное взаимодействие с пользователем. Для этого используются анимации, движения и различные реакции робота.
При этом основное назначение таких устройств заключается в развлечении и создании эффекта взаимодействия с живым существом, а не в мониторинге внимания или поддержании продуктивности пользователя.
Открытые микроконтроллерные платформы (Stack-chan)
Stack-chan — открытая робототехническая платформа на базе микроконтроллера ESP32-S3. Она позволяет изменять как программную, так и аппаратную часть устройства и создавать собственные сценарии поведения.
Однако исходная концепция Stack-chan в большей степени ориентирована на создание диалогового робота-компаньона. Для использования платформы в качестве системы мониторинга внимания потребуется дополнительная реализация компьютерного зрения, анализа состояния пользователя и соответствующих реакций робота.
Программные трекеры времени и продуктивности
Отдельную категорию составляют программные решения для контроля времени работы за компьютером и ограничения времени использования экрана. Такие ­стройство способно перемещаться по рабочей поверхности.
Движение используется как дополнительный способ привлечения внимания пользователя. Например, при обнаружении снижения концентрации робот может выполнить определённое движение или изменить своё положение, чтобы обратить на себя внимание.
2. Гибридная архитектура ESP32-S3 + ПК
Система использует распределение задач между персональным компьютером и микроконтроллером ESP32-S3.
Персональный компьютер выполняет более сложные вычислительные задачи: обработку видеопотока, компьютерное зрение и анализ состояния пользователя. ESP32-S3 отвечает за управление двигателями, дисплеем и другими периферийными устройствами робота.
Такой подход позволяет использовать вычислительные возможности ПК для анализа данных, а микроконтроллер — для быстрого управления аппаратной частью устройства.
3. Настраиваемое поведение робота
Предусматривается возможность настройки поведения устройства в соответствии с предпочтениями пользователя. Можно изменять визуальные реакции дисплея, сценарии движения и условия, при которых робот реагирует на снижение концентрации.
Это позволяет адаптировать устройство под разные сценарии использования и сделать взаимодействие с роботом более удобным.
Таким образом, разрабатываемая система объединяет функции мониторинга внимания, физического взаимодействия и персонализированной обратной связи.
В отличие от обычных программных трекеров продуктивности, система использует физического робота для взаимодействия с пользователем. При этом, в отличие от развлекательных роботов-компаньонов, основное назначение 

\subsection{Подраздел 1.1}

Текст текст текст


\subsection{Подраздел 1.2}

Текст текст текст



\subsection{Подраздел 1.3}


\section{Раздел 2}
Текст текст текст

\subsection{Подраздел 2.1}

Текст текст текст

\subsection{Подраздел 2.2}

Текст текст текст

\subsection{Подраздел 2.3}

Текст текст текст

\section{Раздел 3}
Текст текст текст

\subsection{Подраздел 3.1}

Текст текст текст

\subsection{Подраздел 3.2}

Текст текст текст

\subsection{Подраздел 3.3}

Текст текст текст

\subsection{Подраздел 3.4}

Текст текст текст

\section*{Заключение}
\addcontentsline{toc}{section}{Заключение}
В рамках данного курсового проекта Текст текст текст

\begin{thebibliography}{1}
   \bibitem{s1} Медведев А. В. Оптимизационная система поддержки принятия
решений в бизнес-планировании / Успехи современного естествознания.
– 2015. – №1-4 – С.679-683.

    \bibitem{s3} Лазарев Д. И., Мажан Е. В., Рублев В. В. Особенности процесса принятия решений и анализ проблем управления в современном предприятии / Сборник трудов конференции. – 2016. – С. 212–217.

    \bibitem{s4} Закалин И. Ю. Методы поддержки принятия решений на основе информационных технологий / ВЕСТНИК МАГИСТРАТУРЫ  – 2018. – № 1-1 (76). – С. 23–26.

    \bibitem{s2} Ахрамейко А. А., Хмельницкая И. В. Концептуальные проблемы построения информационно-аналитических систем поддержки принятия решения / Сборник статей. – 2010. – С. 68–72.
\end{thebibliography}

\newpage
\begin{center}
\textbf{Приложение А }

\text{(обязательное)}

\text{Экранные формы}
\end{center}
\addcontentsline{toc}{section}{Приложение А}

\newpage
\begin{center}
\textbf{Приложение Б }

\text{(обязательное)}

\text{Исходный код}
\end{center}
\addcontentsline{toc}{section}{Приложение Б}

\end{document}
